\documentclass[10pt,a4paper,roman]{moderncv}
\moderncvtheme[black]{classic}
\usepackage[utf8]{inputenc}
\usepackage[scale=0.85]{geometry} % La taille pris par le contenu, ici on a 15% de marges.
\usepackage{xcolor}
\usepackage{natbib}
\usepackage{bibentry}
\usepackage[bottom]{footmisc}

\renewcommand*{\namefont}{\fontsize{18}{20}\mdseries\upshape}
\renewcommand*{\titlefont}{\fontsize{14}{20}\mdseries\upshape}

% Define custom type of list
\newcounter{pubctr}
\newenvironment{publist}
  {%
    \begin{list}{[\arabic{pubctr}]}%
      {%
        \setlength{\leftmargin}{1.5em}%
        \setlength{\itemsep}{0.3em}%
        \setlength{\labelsep}{0.5em}%
      }%
    \newcommand{\pubitem}{%
      \refstepcounter{pubctr}%
      \item
    }%
  }
  {%
    \end{list}%
  }

\newcounter{thesisctr}
\newenvironment{thesislist}
  {\begin{list}{[\arabic{thesisctr}]}%
     {\usecounter{thesisctr}
      \setlength{\leftmargin}{1.5em}
      \setlength{\itemsep}{0.3em}
     }}
  {\end{list}}

\nopagenumbers{} % Permet de masquer les numéros de page
 
\title{Research scientist in cloud computing}
\firstname{\textcolor{black}{\textbf{Pierre}}}
\familyname{\textbf{JACQUET}}
\address{1100 rue Notre Dame Ouest,}{H3C 1K3 Montréal, Canada}
\mobile{ +(514) 396-8800}
\email{hidden-mail@domain.com}
\extrainfo{\href{https://jacquetpi.github.io}{https://jacquetpi.github.io}}
 
\begin{document}
\maketitle

 
\section{Work Experience}

\cventry{Jan. 2025 - Dec. 2026}{École de technologie supérieure}{Montréal}{Canada\newline Postdoctoral Researcher\newline Topic: A Better Sharing of Cloud Resources}{In a partnership with OVHcloud Montréal}{}{}
\cventry{Oct. 2021 - Sept. 2024}{Université de Lille, Inria}{Lille}{France\newline Ph.D. Candidate\newline Topic: Enhancing IaaS Consolidation with Resource Oversubscription}{}{}{}
\cventry{Sept. 2018 - Sept. 2021}{Groupe BPCE}{Paris}{France\newline Mainframe Security Architect\newline Apprenticeship. Project management related to the banking security ecosystem}{}{}{}
\cventry{2018}{Direction Générale des Finances Publiques}{Paris}{France\newline Intern. Java developer}{}{}{}

\section{Education}

\cventry{2024}{Université de Lille}{Lille}{France\newline Ph.D. –Computer Science\newline Advisors: Thomas Ledoux, Romain Rouvoy}{}{}
\cventry{2021}{Université de Paris-Est}{Créteil}{France\newline Master of Computer Science}{}{}
 
\section{Publications}
\bibliographystyle{plain}
\nobibliography{cv}

\textbf{International Journal/Conference\footnote{In distributed systems, conferences are often the primary venues for publishing high-impact research and typically involve rigorous peer review. Author order is determined by individual contribution to the paper.} Publications:}

\begin{publist}
{{PUBLICATIONS_MAIN}}
\end{publist}

\newpage
\textbf{Workshop and Short Papers:}

\begin{publist}
{{PUBLICATIONS_WORKSHOP}}
\end{publist}

\textbf{Thesis:}

\begin{thesislist}
{{PUBLICATIONS_THESIS}}
\end{thesislist}

\section{Research artifacts}

This section presents a list of software artifacts for which I directly contributed to the implementation and source code. 
The self-assessment of each artifact follows the criteria defined in the document \href{https://www.inria.fr/sites/default/files/2021-01/Criteria%20software%20self%20assessment.pdf}{Criteria for Software Self-Assessment v3 published by Inria}.\\

\begin{itemize}
{{ARTIFACTS_SECTION}}
\end{itemize}
 
\section{Teaching experience}
 
\textbf{Université de Lille}\\

During my PhD, I held two part-time teaching contracts at the University of Lille over the 2022--23 and 2023--24 academic years, for a total of 102 hours of teaching (equivalent TD). 
My responsibilities consisted in delivering tutorials and practical sessions, grading, and supervising students, based on existing pedagogical material and course content.\\
 
\begin{itemize}
    \item Introduction to Object-oriented programming. Fall22, Fall23 (Teaching Assistant, 30 students, L2)
    \item Advanced Object-oriented programming. Spring23 (Teaching Assistant, 20 students, M1)
    \item Introduction to Software security. Spring23 (Teaching Assistant, 20 students, M1)
\end{itemize}

\hfill \break
\textbf{CentraleSupélec}\\

During my postdoctoral fellowship, I designed and supervised a 35-hour student project on the environmental impact of cloud computing, directly grounded in my research. 
This included the definition of learning objectives, project scope, creation of pedagogical content (software tools and datasets), as well as evaluation criteria and grading.\\

\begin{itemize}
    \item Environmental Impact of Cloud Computing (Hackathon). Spring26 (Consultant, 30 students, ING1)
\end{itemize}

% \section{Supervision}

\section{Research Grants} % Grants and fellowships

\begin{itemize}
    \item {\textbf{Mitacs Elevate Fellowship}, Canada, 2024 — \$120,000 CAD for postdoctoral research on cloud resource optimization in collaboration with OVHcloud.}
\end{itemize}

\section{Service}

{{SERVICE_SECTION}}

\section{Distinctions:}

During my career, I have received the following awards and distinctions:\\

Travel grant, \href{https://2026.eurosys.org/grants.html}{EuroSys 2026} — awarded to present on main track\\
24-month scholarship, ÉTS Montréal \href{https://www.mitacs.ca/our-programs/elevate/}{(MITACS Elevate program)} — awarded after an application to an excellence fellowship\\
Best Contribution Award, \href{https://community.ibm.com/community/user/blogs/meredith-stowell1/2021/03/31/announcing-the-2020-master-the-mainframe-winners}{IBM Master the Mainframe} (2020) — international innovation competition with 25{,}000 participants\\
Jury Prize, \href{https://www.agorize.com/en/challenges/black-out-2/pages/saison-1}{Safran Black-Out Challenge} (2019) — national innovation competition with 650 participants\\

\section{Outreach}

\textbf{Talks:}\\ % Invited talks

I was invited to present my work in the following events:\\

\begin{itemize}
\setlength{\itemsep}{0.35em}
{{OUTREACH_TALKS}}
\end{itemize}

\hfill \break
\textbf{Interviews:}\newline

{{OUTREACH_INTERVIEWS}}

\hfill \break
\textbf{Press Articles:}\newline

{{OUTREACH_PRESS}}

\hfill \break
\textbf{Industrial transferts:}\\

\textit{PhD -- Inria Défi project with OVHcloud}\\

My PhD was conducted within an Inria \textit{Défi} project involving OVHcloud as an industrial partner. The collaboration included access to internal datasets and infrastructure metrics, joint definition of research questions, regular technical exchanges, and transfer of research results related to resource management.\\
Contract type: Inria Défi collaborative framework. \\
Participants: Inria, OVHcloud.\\
Duration: 3 years.\\

\textit{Postdoc -- MITACS industrial partnership with OVHcloud}\\

My postdoctoral fellowship was hosted at OVHcloud Montréal under a two-year MITACS contract involving OVHcloud and ÉTS Montréal. The objectives included the integration of research results on orchestration (completed in early 2026), advanced GPU orchestration, and the improvement of carbon accounting methodologies for cloud infrastructures.\\  
Contract type: MITACS industrial research contract.\\ 
Participants: ÉTS Montreal, OVHcloud, MITACS.\\ 
Duration: 2 years.

\section{References}

\begin{tabular}{lll}
\begin{minipage}[t]{2.5in}
\textbf{Prof.\ Camille Coti}\\
1100, rue Notre-Dame Ouest\\
Montréal (Québec) H3C 1K3, CA\\
Phone: +1 514 396-8800\\
\href{mailto:hidden-mail@domain.com}{hidden-mail@domain.com}
\end{minipage} &
\begin{minipage}[t]{2.5in}
\textbf{Prof.\ Marcos Dias De Assunção}\\
1100, rue Notre-Dame Ouest\\
Montréal (Québec) H3C 1K3, CA\\
Phone: +1 514 396-8800\\
\href{mailto:hidden-mail@domain.com}{hidden-mail@domain.com}
\end{minipage} &
\begin{minipage}[t]{2.5in}
\textbf{Prof.\ Romain Rouvoy}\\
40, avenue Halley - Bât A\\
59650 Villeneuve d'Ascq, France\\
Phone: +33 0 00 00 00 00\\
\href{mailto:hidden-mail@domain.com}{hidden-mail@domain.com}
\end{minipage} \\
\\
\begin{minipage}[t]{2.5in}
\textbf{Prof.\ Thomas Ledoux}\\
4 rue Alfred Kastler\\
44307 Nantes, France\\
Phone: +33 0 00 00 00 00\\
\href{mailto:hidden-mail@domain.com}{hidden-mail@domain.com}
\end{minipage} &
\end{tabular}

\end{document}
